\begin{enumerate}

\item El desplazamiento normalizado se prescribe como promedio de pendientes.

\[\fitmath{\mathrm{dd}=\frac{\mathrm{me}+\mathrm{ms}}{2}}\]

\[\fitmath{\mathrm{dd}=\frac{(0.1278)+(1.38)}{2}\approx 0.7539\,}\]

\item El calado de cuerda es atan d.

\[\fitmath{\mathrm{chi}=\arctan\left(\frac{\mathrm{me}+\mathrm{ms}}{2}\right)}\]

\[\fitmath{\mathrm{chi}=\arctan\left(\frac{(0.1278)+(1.38)}{2}\right)\approx 0.6459924\,\mathrm{rad}}\]

\item La cuerda real se deduce del desplazamiento final.

\[\fitmath{\mathrm{chord}=\mathrm{cxg}\cdot \sqrt{1+{\left(\frac{\mathrm{me}+\mathrm{ms}}{2}\right)}^{2}}}\]

\[\fitmath{\mathrm{chord}=(0.046)\cdot \sqrt{1+{\left(\frac{(0.1278)+(1.38)}{2}\right)}^{2}}\approx 0.05760782\,\mathrm{m}}\]

\end{enumerate}
